%!TEX root = ../../main.tex
%% 1.4 기여와 비주장 — 근거: journal 06_INTRODUCTION_BRIEF P7, JOURNAL_FINISH_LOCK Forbidden, 01_MRQ_CLAIM_MAP §3
\section{기여와 비주장}
\label{sec:intro-contrib}

본 논문의 기여는 다음 네 가지이다.

\begin{enumerate}[leftmargin=2em]
  \item \textbf{근거 있는 교전 사례 구성 (M-RQ1).} 시간 경계 $G$와 공간 경계 $D$를 코퍼스에서 추정하고, 게이트 상수 $R$, $B$를 게임 규칙에 고정하며, 교전 집합을 결정하는 모든 상수를 출처 분류와 함께 표로 남긴다. 규모 계급을 작은 쪽 참여 인원으로 정의하고, 한타 T와 소규모 교전 S를 별개의 코호트로 둔다.
  \item \textbf{경기 승패와 연결된 결과 가치 (M-RQ2).} 경기 승패로 학습해 동결한 승률 평가기 $\Vhat$로 교전 구간 전후의 추정 승률 변화 방향 $\Ysvi = \ind[\dV>0]$을 정의하고, 평가기 품질, 변화의 방향·평균·규모의 분리, 비교전 구간 대조, 물질적 결과와 다음 오브젝트와의 대응, 지평 안정성으로 이 측정이 무엇을 재는지 한정한다.
  \item \textbf{초기 우세와 시간을 넘는 예측 가능성 (M-RQ3).} 코호트별로 교전 전 상태의 정규화 로지스틱 모델 $q$가 승률·시간의 스플라인 기준선 $\PTflex$보다 Brier 점수를 낮추는지 동일 행의 짝 대비와 경기 단위 부트스트랩 구간으로 답한다. 한타 모델의 소규모 교전 전이와 합동 학습의 효과를 예비 선언된 이차 대비로 보고한다.
  \item \textbf{의미와 범위의 검증 (M-RQ4).} 균형 구간, 작은 변화 제외, 점수 분해, 이후 패치·다른 지역의 점수 전용 평가로 이득이 어디까지 유지되는지 보고한다.
\end{enumerate}

동시에 본 논문은 다음을 \emph{주장하지 않는다}. 이 목록은 동반 저널 원고의 마감 잠금과 규모분리 계약에서 금지 항목으로 고정된 것이며, 제 \ref{ch:discussion} 장에서 표로 다시 확인한다.

\begin{itemize}[leftmargin=1.5em]
  \item 공개 정보로는 어떤 설계에서도 특정 AUC를 넘을 수 없다는 상한 주장.
  \item 균형 구간의 약한 이득이 관측되지 않은 전투 숙련의 증거라는 주장.
  \item 재보정만으로 외부 전이가 복구된다거나 외부 실패가 보정만의 문제라는 주장.
  \item 최적 아키텍처나 망라적 모델 순위.
  \item 물질적 결과·다음 오브젝트와의 대응이 $q$의 정확도이거나 인과적 처치 효과라는 주장.
  \item 한타 또는 소규모 교전이 ``더 예측 가능하다''는 주장과 규모 기울기 일반.
  \item T와 S의 차이에 대한 기전 설명.
  \item 합동 T$\cup$S, pick, 소규모 외부 코호트의 결과를 해당 코호트의 주 대비 대신 쓰는 것.
\end{itemize}
